\clearpage
\renewcommand{\BookRunningTitle}{Geometrical Aspect of Ternary Forms}
\SourceLeaf{315}
\pdfbookmark[0]{Geometrical Aspect of Ternary Forms}{ternary-geometry}
\begin{center}\Large\bfseries GEOMETRICAL ASPECT OF TERNARY FORMS.\end{center}
\BookRule{30mm}

Let a point in space $(0)$ be chosen as the origin. Let the displacements from it to three other points $P,P',P''$, which do not lie in one plane with it, be $t,t',t''$, respectively; whenever confusion is impossible, the points $P,P',P''$ themselves may be denoted by $(t)$. $(t')$, $(t'')$.

Further, let $(t,t')$ generally denote the product of the lengths of the two lines $t,t'$ and the cosine of their inclination, etc.

In general, $(\alpha t+\alpha't'+\alpha''t''+\ldots\quad\beta u+\beta'u'+\beta''u''+\mathinner{.\,.})$ is obtained by carrying out the multiplication
\[
(\alpha t+\alpha't'+\alpha''t''+\mathinner{.\,.})\times(\beta u+\beta'u'+\beta''u''+\mathinner{.\,.})
\]
and replacing $tu,\allowbreak tu',\allowbreak tu'',\allowbreak t'u',\allowbreak t'u''\mathinner{.\,.}$ and so on by $(t,u),\allowbreak(t,u'),\allowbreak(t,u''),\allowbreak(t',u'),\allowbreak(t',u'')$ and so on.

Every point in space can be represented by a trinomial
\[
(xt+x't'+x''t'')
\]

For all points in a given plane, an equation
\[
\lambda x+\lambda'x'+\lambda''x''=L
\]
\SourceLeaf{316}
will then hold, where $\lambda,\lambda',\lambda'',L$ are fixed numbers. For a plane through the three points $\mu t,\mu't',\mu''t''$,
\[
\lambda\mu=\lambda'\mu'=\lambda''\mu''=L
\]

Writing
\[
\begin{gathered}
(t,t)=a,\quad(t',t')=a',\quad(t'',t'')=a'',\\
(t',t'')=b,\quad(t,t'')=b',\quad(t,t')=b''
\end{gathered}
\]
and
\[
\begin{gathered}
a'a''-bb=A,\quad aa''-b'b'=A',\quad aa'-b''b''=A''\\
b'b''-ab=B,\quad bb''-a'b'=B',\quad bb'-a''b''=B''\\
D=aa'a''+2bb'b''-abb-a'b'b'-a''b''b''
\end{gathered}
\]
gives
\[
\begin{array}{rclll}
T&=&At+B''t'+B't''&\text{perpendicular to}&t'\text{ and }t''\\
T'&=&B''t+A't'+Bt''&&t\text{ and }t''\\
T''&=&B't+Bt'+A''t''&&t\text{ and }t'
\end{array}
\]
and, in general, if
\begingroup\postdisplaypenalty=10000
\[
\lambda x+\lambda'x'+\lambda''x''=L
\]\endgroup
is the equation of a plane, the line
\[
\lambda T+\lambda'T'+\lambda''T''
\]
will be perpendicular to that plane.

Further,
\[
\begin{aligned}
aT+b''T'+b'T''&=Dt\\
b''T+a'T'+bT''&=Dt'\\
b'T+bT'+a''T''&=Dt''
\end{aligned}
\]
and the lines $t,t',t''$ are perpendicular to the planes whose equations
\[
\begin{aligned}
ax+b''x'+b'x''&=\mathrm{Const}\\
b''x+a'x'+bx''&=\mathrm{Const}\\
b'x+bx'+a''x''&=\mathrm{Const}
\end{aligned}
\]

\SourceLeaf{317}
Twice the area of the triangle through the points $mt,m't',m''t''$ equals the square root of the value of the form
\[
F\ldots\left(\begin{smallmatrix}A,&A',&A''\\B,&B',&B''\end{smallmatrix}\right)
\]
on substituting $X=m'm'',X'=mm'',X''=mm'$, while six times the volume of the pyramid formed with the point $0$ becomes $=mm'm''\surd D$; hence the altitude is
\[
=\surd\frac{D}{F(\frac1m,\frac1{m'},\frac1{m''})}
\]
$T,T',T''$ relate to the form $\left(\begin{smallmatrix}AD,&A'D,&A''D\\BD,&B'D,&B''D\end{smallmatrix}\right)$ as $t,t',t''$ do to $\left(\begin{smallmatrix}a,&a',&a''\\b,&b',&b''\end{smallmatrix}\right)$

The three roots of the equation
\[
0=p^3-pp(a+a'+a'')+p(A+A'+A'')-D
\]
give the squares of the three principal axes of an ellipsoid inscribed in the parallelepiped to which the positive ternary form
\[
\left(\begin{smallmatrix}a,&a',&a''\\b,&b',&b''\end{smallmatrix}\right)
\quad\text{with adjoint }\left(\begin{smallmatrix}A,&A',&A''\\B,&B',&B''\end{smallmatrix}\right)
\quad\text{and det. }=-D
\]
relates.
\BookRule{.30\textwidth}

\subsection*{Expressing spatial relations relative to a given tetrahedron.}
\pdfbookmark[1]{Spatial relations relative to a given tetrahedron}{ternary-tetraeder}
Let $(0),(1),(2),(3)$ denote the four vertices, opposite faces and altitudes. Each point in space $P$, relative to an arbitrary origin $M$, then has four coordinates $x,x',x'',x'''$, subject, however, to the relation
\[
x+x'+x''+x'''=0
\]
Here $x$ denotes the quotient obtained by dividing the distance of the
\SourceLeaf{318}
point $P$ from a plane drawn through $M$ parallel to the plane $(0)$ by the altitude $(0)$, and so on.

In general,
\[
\begin{aligned}
-(PM)^2={}&xx'(01)^2+xx''(02)^2+xx'''(03)^2\\
&+x'x''(12)^2+x'x'''(13)^2+x''x'''(23)^2
\end{aligned}
\]
\BookRule{.30\textwidth}

\emph{The fundamental law of crystallization} can be stated most briefly as follows: any five planes occurring in it satisfy the following relation:

If their normals on the sphere are $(0),(1),(2),(3),(4)$, the products $\sin102\pdot\sin304$, $\sin103\pdot\sin204$, $\sin203\pdot\sin104$ are always in a rational ratio; if this is $\alpha:\beta:\gamma$, then $\beta=\alpha+\gamma$.

If the coordinates of the $5$ points on the sphere are
\[
\left.\begin{array}{ccc}
a&b&c\\a'&b'&c'\\a''&b''&c''\\a'''&b'''&c'''\\0&0&1
\end{array}\right\}
\quad\text{then the products must}\quad
\begin{gathered}
(ab'-ba')\pdot(a''b'''-b''a''')\\
(ab''-ba'')\pdot(a'''b'-b'''a')\\
(ab'''-ba''')\pdot(a'b''-b'a'')
\end{gathered}
\]
be in a rational ratio.

In general, let $1,2,3,4,5$ be the $5$ points on the sphere and $0$ its centre; then, if $12$ denotes the volume of the tetrahedron $0345$,
\[
\left.\begin{array}{c}23\pdot45\\24\pdot53\\25\pdot34\end{array}\right\}
\quad\text{are in a rational ratio}
\]
and likewise
\[
\left.\begin{array}{c}12\pdot34\\13\pdot42\\14\pdot23\end{array}\right\}\quad\text{and so on}
\]
